% ================================
% PROBLEM ARTICULATION
% ================================
\section{Problem Articulation}

At the continuum level, electrical conduction in a deformable material can be described 
by a conductivity equation in which the electrical conductivity depends on the local mechanical strain. 
Such models are widely used in engineering applications and can be solved efficiently using numerical methods. 
However, the central quantity in this description, namely the strain-dependent conductivity $\sigma(\varepsilon)$, 
is often introduced phenomenologically, without an explicit connection to the underlying material physics.

From a materials physics perspective, this raises a fundamental question: 
how does mechanical deformation at the atomic scale translate into a change in macroscopic electrical conductivity?

In semiconductors, strain modifies the crystal structure and therefore the electronic band structure of the material. 
This leads to changes in quantities such as band energies, effective masses, and carrier velocities. 
As a consequence, the transport properties of charge carriers are altered, affecting mobility and electrical conductivity. 
These mechanisms are described in semiconductor theory through band-structure analysis and transport models 
such as the semiclassical Boltzmann transport equation \cite{hamaguchi}.

This observation suggests a natural multiscale structure of the problem. 
At the microscopic level, strain influences the electronic structure of the material. 
At the transport level, the modified electronic structure determines carrier dynamics and effective conductivity. 
At the continuum level, this conductivity enters as a constitutive parameter in the electrical conduction equation.

The core challenge of this work is therefore to connect these levels in a computationally consistent way. 
Rather than treating $\sigma(\varepsilon)$ purely as a fitting function, the goal is to interpret it as an effective quantity 
that reflects underlying electronic and transport mechanisms, while remaining suitable for use in continuum simulations.

This perspective allows the continuum model to serve not only as a numerical tool, but also as a bridge between 
microscopic physical processes and macroscopic device behavior.

In particular, this framework enables the incorporation of physically informed or data-driven constitutive relations, depending on the level of detail available for the underlying material.